section{Insights from Work Undertaken} \label{section:insights_from_work}

\subsection{What was learned from each system}

\subsubsection{Carbon (C) – data incompleteness and implications}

The carbon dataset remains incomplete and cannot yet support definitive conclusions. Preliminary indications suggest that strong bonding and limited configurational diversity may pose challenges analogous to those observed in Si.

\subsubsection{Silicon (Si) – persistent poor rank correlation}

Across all iterations, Si consistently exhibited poor MLP--DFT agreement. This suggests intrinsic modelling difficulties rather than statistical noise. The failures appear systematic, reinforcing the need for targeted improvements.

\subsubsection{Germanium (Ge) and Tin (Sn) – favourable MLP--DFT agreement}

For both Ge and Sn, MLPs reliably preserved DFT-derived rankings. These systems are likely well-represented in the MACE training dataset, and their more metallic bonding character may aid accurate regression of energies.

\subsection{Contrast perfect alloys vs heterostructures}

Perfect alloy systems showed tighter rank correlation and lower error dispersion than heterostructures.

\begin{table}[h]
    \centering
    \caption{Comparison of MLP performance on perfect alloys and heterostructures.}
    \begin{tabular}{lccc}
        System & Spearman & Kendall & Top-10 Overlap \\
        \hline
        Perfect Alloys & 1.00 & 1.00 & 6 \\
        Ge|Si Heterostructure & 0.73 & 0.55 & 6 \\
    \end{tabular}
\end{table}

\section{Challenges in Modelling Silicon Interfaces} \label{section:challenges_in_modelling_silicon_interfaces}

\subsection{Discussion of why MLPs fail for Si:}

\subsubsection{Local bonding features?}

MLPs may struggle with the directional covalent bonding and short bond lengths characteristic of Si. These features introduce sensitivity to structural distortions that are difficult to learn from sparse data.

\subsubsection{Underrepresented environments in training?}

It is likely that the MACE training set lacks sufficient coverage of interfacial Si environments, especially those with reconstructed surfaces and strained bonds.

\subsection{Open questions}

\subsubsection{Is Si anomalous, or are MLP limitations broader?}

The current data cannot confirm whether Si represents a broader warning about MLP application or a special case. Additional experiments with other covalently bonded systems are required to determine generalisability.